% Компілювати: xelatex Lecture_12_Instrumental_Analysis_Math.tex
\documentclass[11pt,a4paper]{article}
\usepackage{fontspec}
\usepackage{polyglossia}
\setmainlanguage{ukrainian}
\setmainfont{DejaVu Serif}
\usepackage{amsmath,amssymb}
\usepackage[margin=22mm]{geometry}
\setlength{\parindent}{0pt}
\setlength{\parskip}{5pt plus 2pt}
\setlength{\emergencystretch}{3em}
\raggedbottom
\newcommand{\unit}[1]{\,\mathrm{#1}}
\newcommand{\eqv}{\mathrm{eq}}
\newcommand{\Vendpoint}{V_{\mathrm{end}}}
\title{Лекція 12. Титрування, інструментальний аналіз і вимірювальні системи\\[4pt]\large Математичний конспект}
\author{}
\date{}
\begin{document}
\raggedright
\maketitle

\section{Титриметричний аналіз}

\textbf{Стехіометричний баланс.}
Нехай аналіт \(a\) і титрант \(t\) реагують за співвідношенням
\(\nu_a a+\nu_t t\to\mathrm{products}\), де \(\nu_a,\nu_t>0\) ---
безрозмірні стехіометричні коефіцієнти. За повного перебігу реакції
без побічного споживання титранту в точці еквівалентності
\[
  \frac{c_a V_a}{\nu_a}=\frac{c_t V_{\eqv}}{\nu_t},
  \qquad
  c_a=\frac{\nu_a}{\nu_t}\frac{c_t V_{\eqv}}{V_a},
  \qquad
  V_{\eqv}=\frac{\nu_t}{\nu_a}\frac{c_a V_a}{c_t}.
\]
\(c_a,c_t\) --- молярні концентрації в \(\mathrm{mol/L}\),
\(V_a>0\) --- об'єм аліквоти, \(V_{\eqv}\) --- сукупний об'єм
титранту в точці еквівалентності. У балансі кількостей речовини
обидва об'єми підставляють у літрах; у відношенні об'ємів
можна використовувати мілілітри для обох. Для \(\nu_a=\nu_t\),
\(V_a=20.0\unit{mL}\), \(c_t=0.100\unit{mol/L}\) та
\(V_{\eqv}=8.00\unit{mL}\) маємо \(c_a=0.0400\unit{mol/L}\).

\textbf{Кінцева точка за дискретними відліками.}
Для об'ємів \(V_0<V_1<\dots<V_N\), у \(\mathrm{mL}\), і
відповідей \(y_i\) (наприклад, безрозмірного pH) визначимо
\[
  D_i=\frac{y_{i+1}-y_i}{V_{i+1}-V_i},
  \qquad
  i_* \in \mathop{\mathrm{arg\,max}}_{0\le i<N}|D_i|,
  \qquad
  \Vendpoint=\frac{V_{i_*}+V_{i_*+1}}{2}.
\]
Для pH величина \(D_i\) має одиницю \(\mathrm{mL^{-1}}\).
Модуль допускає як зростання, так і спадання сигналу.
За кількох максимумів та без додаткового правила \(\Vendpoint\)
не визначено однозначно; максимум на краю ряду не дає даних
з обох боків переходу. Метод придатний для одного
виразного, достатньо густо відліченого переходу.
Рівність \(\Vendpoint=V_{\eqv}\) не гарантується.

За фіксованих \(c_t,V_a,\nu_a,\nu_t\), підстановка
\(\Vendpoint\) замість \(V_{\eqv}\) змінює оцінку концентрації на
\[
  \Delta c_a =
  \frac{\nu_a c_t}{\nu_t V_a}\bigl(\Vendpoint-V_{\eqv}\bigr).
\]
Лише за умови \(V_{\eqv}\in[V_{i_*},V_{i_*+1}]\) похибка
об'єму за модулем не перевищує половини ширини цього інтервалу;
ця межа не є повною невизначеністю аналізу.

\section{Мас-спектрометричний аналіз}

\textbf{Шкала \(m/z\).}
Нехай маса іона \(m_i\) виражена в \(\mathrm{kg}\),
\(m_u\) --- атомна одиниця маси (\(\mathrm{kg}\)),
\(q_i=z_i e\) --- заряд (\(\mathrm{C}\)), \(z_i\in\mathbb Z\setminus\{0\}\)
--- зарядове число, \(e\) --- елементарний заряд (\(\mathrm{C}\)).
Координата на числовій шкалі мас-спектра:
\[
  r_i=\frac{m_i/m_u}{|z_i|}.
\]
\(r_i\) не має розмірності: чисельник є числовим значенням
маси в дальтонах. Для номінального \(m/z\) замість \(m_i/m_u\)
використовують суму масових чисел атомів за прийнятою ізотопною
моделлю. Фізичне відношення \(m_i/|q_i|\), на відміну від \(r_i\),
має одиницю \(\mathrm{kg/C}\).
Іони з різними масами та зарядами можуть мати однаковий \(r_i\).

\textbf{Відбір кандидатів.}
Для піків \((r_j,I_j)\) одного спектра нехай
\(I_{\max}=\max_j I_j>0\) і \(J_j=100 I_j/I_{\max}\), у відсотках;
\(I_j\ge0\) та \(I_{\max}\) виміряно в однакових одиницях.
Кандидат \(s\) із маркерами \(M_s=\{\mu_1,\ldots,\mu_k\}\)
проходить поріг, якщо існує ін'єктивне зіставлення \(\pi\)
маркерів з виміряними піками таке, що для кожного маркера
\[
   |r_{\pi(l)}-\mu_l|\le \delta,\qquad
   J_{\pi(l)}\ge J_{\min},\qquad l=1,\ldots,k.
\]
\(\delta\ge0\) --- допуск у числовій шкалі \(m/z\);
\(J_{\min}\in[0,100]\) --- поріг відносної інтенсивності.
Умова однакового типу іонізації та сумісних параметрів
спектрів входить у припущення зіставлення.
Нуль кандидатів означає відсутність збігу за критерієм,
один --- єдиний збіг у заданій бібліотеці,
кілька --- неоднозначність.

\textbf{Кількісна оцінка з внутрішнім стандартом.}
Для сигналів \(S_a\) та \(S_{\mathrm{is}}>0\) в однакових одиницях
можна застосувати калібрувальну модель
\[
  R=\frac{S_a}{S_{\mathrm{is}}}=b+kc,
  \qquad c=\frac{R-b}{k}.
\]
\(R,b\) безрозмірні; для \(c\) у \(\mathrm{mg/L}\)
коефіцієнт \(k\neq0\) має одиницю \(\mathrm{L/mg}\).
Модель потребує валідованого діапазону, сталих умов
введення внутрішнього стандарту та роздільно вимірюваних сигналів.

\section{Хроматографічний аналіз}

\textbf{Час утримування та площа.}
Для часу введення проби \(t_{\rm inj}\) та часу максимуму піку
\(t_{\rm max}\) час утримування \(t_R=t_{\rm max}-t_{\rm inj}\)
вимірюється в секундах, як і обидві часові мітки.
Розглядаємо розділений пік у вікні \([t_a,t_b]\),
спостереження \(s(t)\) і оцінку базової лінії \(b(t)\),
обидві в \(\mathrm{mV}\). Чиста площа
\[
    P=\int_{t_a}^{t_b}(s(t)-b(t))\,dt
\]
вимірюється в \(\mathrm{mV\,s}\).
Для відліків \(t_0<\dots<t_N\) у секундах та меж вікна, що
збігаються з вузлами \(t_l,t_r\), складена формула трапецій:
\[
   P_{\mathrm{tr}} =
   \sum_{i=l}^{r-1}\frac{u_i+u_{i+1}}{2}(t_{i+1}-t_i),
   \qquad u_i=s(t_i)-b(t_i).
\]
Формула точна для кусково-лінійного \(u\) між вузлами.
Якщо \(u\in C^2([t_l,t_r])\) і \(\max|u''|\le M\)
(\(\mathrm{mV/s^2}\)), то
\[
   |P-P_{\mathrm{tr}}|
     \le \frac{M}{12}\sum_{i=l}^{r-1}(t_{i+1}-t_i)^3
     \le \frac{M}{12}(t_r-t_l)\max_i(t_{i+1}-t_i)^2 ,
\]
де максимум кроку береться у вікні. Ця межа стосується
квадратурної похибки, а не похибки визначення базової лінії.
Для \(u=(0,6,0)\unit{mV}\) у моменти \(t=(0,1,3)\unit{s}\)
площа становить \(9\unit{mV\,s}\).

\textbf{Обернене калібрування за площею.}
За фіксованих умов пробовведення, поділу й детектування можна
використати \(P=P_0+kc\) для конкретного аналіту;
тоді \(c=(P-P_0)/k\). Для \(c\) у \(\mathrm{mg/L}\):
\[
   [P_0]=\mathrm{mV\,s},\qquad
   [k]=\mathrm{mV\,s\,L/mg},\qquad k>0 .
\]
Наприклад, \(P_0=1\unit{mV\,s}\) та \(k=4\unit{mV\,s\,L/mg}\)
дають \(c=2\unit{mg/L}\) для \(P=9\unit{mV\,s}\).
Обернену модель застосовують лише для площ та концентрацій
у встановленому калібрувальному діапазоні.
Інтеграл перекритих піків не є площею окремого компонента
без додаткового способу розділення внесків.

\section{Промисловий аналізатор та автоматизована система}

\textbf{Номінальна затримка.}
Для ділянки з внутрішнім об'ємом \(V>0\) у \(\mathrm{mL}\)
і сталою об'ємною витратою \(q>0\) у \(\mathrm{mL/s}\)
модель перенесення без змішування дає
\[
   \tau_{\rm line}=\frac{V}{q},\qquad
   \tau_{\rm total}=\tau_{\rm line}+
                      \tau_{\rm prep}+\tau_{\rm instr}+\tau_{\rm data}.
\]
Усі \(\tau\) вимірюються в секундах, доданки відповідають
послідовним етапам без перекриття часу.
Для \(V=90\unit{mL}\), \(q=3\unit{mL/s}\) та 12~с решти
етапів \(\tau_{\rm total}=42\unit{s}\).
Обернене відношення \(q=V/\tau_{\rm line}\) повертає
\(3\unit{mL/s}\) при \(\tau_{\rm line}=30\unit{s}\).
Модель не описує дисперсію часу перебування, адсорбцію,
реакції у лінії чи розмивання сходинки концентрації.

\textbf{Придатність специфікації та окремого відліку.}
Нехай \(D_{\rm req}=[c_L,c_H]\) --- потрібний діапазон,
\(D_{\rm cal}=[\ell,h]\) --- калібрований діапазон,
усі межі в \(\mathrm{mg/L}\) і впорядковані.
Для встановленого обмеження \(T_{\rm lim}\ge0\) у секундах
умови придатності специфікації:
\[
   D_{\rm req}\subseteq D_{\rm cal}
   \ \Longleftrightarrow\ 
   \ell\le c_L\le c_H\le h,
   \qquad
   \tau_{\rm total}\le T_{\rm lim}.
\]
Для окремого запису з оцінкою \(\hat c\) у \(\mathrm{mg/L}\)
та сумісними мітками \(t_{\rm samp},t_{\rm read}\) у секундах
вік даних \(a=t_{\rm read}-t_{\rm samp}\).
За навчальною моделлю запис прийнятний, якщо вимірювальна
процедура повернула справний статус, \(\hat c\in D_{\rm cal}\)
та \(0\le a\le T_{\rm lim}\).
Порушення діапазону й прострочення можуть існувати одночасно;
такі умови перевіряють незалежно. Вік конкретного запису,
на відміну від \(\tau_{\rm total}\), включає й час до фактичного
читання результату.

\end{document}
